Finding optimal transport map components requires navigating a high-dimensional polynomial coefficient space to minimise the empirical log-likelihood subject to polyhedral monotonicity constraints. To solve this constrained problem efficiently across large particle sets, we introduce a highly scalable, inexact Projected Gradient Descent (PGD) architecture. The framework integrates stochastic gradient evaluation to navigate complex objective landscapes, iterative hard thresholding to enforce structural sparsity, and dynamic projection scheduling to circumvent the prohibitive computational costs of exact Euclidean projections.

\section*{Notation}
We denote set cardinality as $|\cdot|$. The Hadamard product, denoted by the operator $\odot$, represents the element-wise multiplication of two vectors or matrices of identical dimensions.
The outer learning PGD loop is indexed by $t \in \{0, 1, \dots, T-1\}$, where $T-1$ is the outer loop iteration budget, while the inner Dykstra projection cycle is indexed by $m$ and constrained by a dynamic iteration ceiling $M^{(t)}$, which scales according to a base budget $M_{\mathrm{base}}$ and an inexact growth exponent $p$. The stochastic gradient $g^{(t)}$ is evaluated over a random subset of particle indices $\mathcal{B}^{(t)} \subset \{1, 2, \dots, n\}$ with cardinality $B$. The learning rate is denoted as $\eta^{(t)}$, and decays according to an initial rate $\eta^\circ$ and a decay parameter $\gamma$. For the iterative hard thresholding procedure, we define a boolean active mask vector $\mathcal{M}^{(t)}$.

